\subsection{Metodología}

El diseño de la red hipotética para el edificio de salones (A) se gestionará utilizando la metodología del ciclo de vida de la red PPDIOO (Preparar, Planificar, Diseñar, Implementar, Operar y Optimizar). Este marco estructurado garantiza el cumplimiento de los requisitos técnicos y operativos del proyecto.

El objetivo principal es usar este modelo para ejecutar de manera ordenada las seis fases de entrega definidas en el documento del proyecto. La correspondencia entre PPDIOO y las fases del proyecto será la siguiente:

\textbf{Fase 1}. Planeación y diseño. Esta fase se cubre con las tres primeras etapas de PPDIOO:

Preparar (Prepare): Se definen los objetivos del proyecto. En nuestro caso: diseñar la red del Edificio A con "Bajo presupuesto" (\$490,200) , asegurar conectividad con el Edificio de Innovación , implementar CCTV , VoIP y una red inalámbrica "muy estable".

Planificar (Plan): Se identifican los recursos, se crea el cronograma de actividades y se detallan los costos de material y mano de obra para la cotización, asegurando no exceder el presupuesto.

Diseñar (Design): Se crea la solución técnica detallada, incluyendo la topología de red, selección de equipo activo y pasivo (Cat 6 y Fibra OM3), el diagrama de red, la tabla de direccionamiento IP y el diseño de la WLAN segmentada (alumnos/docentes).

\textbf{Fase 2}. Implementación de la red (configuraciones lógicas). Corresponde a la etapa Implementar (Implement) de PPDIOO. En este punto se realizan las configuraciones lógicas de switches, routers y puntos de acceso conforme al diseño.

\textbf{Fase 3}. Monitoreo y análisis de la red. Corresponde a la etapa Operar (Operate): verificar que la red y los servicios (VoIP, CCTV, Wi\-/Fi y enlaces) funcionen correctamente. Se despliegan herramientas de monitoreo, se configuran dashboards y alertas, y se generan reportes periódicos que confirman la operación y permiten detectar y corregir anomalías a tiempo.

\textbf{Fase 4}. Implementación de seguridad. Ambas corresponden a la etapa Implementar (Implement). Aquí se configuran las políticas de seguridad: segmentación de WLAN, creación de VLAN (alumnos, docentes, VoIP, CCTV) y listas de control de acceso para proteger la red.

\textbf{Fase 5}. Administración de fallas. Se engloba en la etapa Operar (Operate) y se enfoca en la gestión diaria de la red, lo que incluye:

El monitoreo constante para establecer un rendimiento base y detectar anomalías (Monitoreo y análisis).

La detección, diagnóstico y resolución proactiva de problemas (Administración de fallas) para mantener la red "muy estable".

\textbf{Fase 6}. Calidad en el servicio. Esta fase corresponde a la etapa Optimizar (Optimize). Con la red operativa y monitoreada, se aplican políticas de calidad de servicio para priorizar el tráfico sensible como VoIP y video de CCTV.

El uso de PPDIOO como metodología de trabajo garantiza que cada entregable del proyecto sea un resultado lógico del ciclo de vida de la red.

